\documentclass{uai2026} % for initial submission
%\documentclass[accepted]{uai2026}

\usepackage[american]{babel}
\usepackage{natbib}
  \bibliographystyle{plainnat}
  \renewcommand{\bibsection}{\subsubsection*{References}}
\usepackage{mathtools}
\usepackage{amssymb}
\usepackage{amsthm}
\usepackage{booktabs}
\usepackage{xcolor}
\usepackage{hyperref}

% Optional: use existing macro collections from the template repo.
\input{math_commands}
\input{defns}

% === Draft controls ===
\newif\ifdraft
\drafttrue % set to \draftfalse before submission

\newcommand{\todo}[1]{\textcolor{red}{[TODO: #1]}}
\newcommand{\plan}[1]{%
  \ifdraft
    \paragraph{Draft plan.}#1
  \fi
}

% === Theorem environments (minimal) ===
\theoremstyle{plain}
\newtheorem{theorem}{Theorem}
\newtheorem{lemma}{Lemma}
\newtheorem{proposition}{Proposition}
\theoremstyle{definition}
\newtheorem{definition}{Definition}

% === Notation shortcuts ===
\newcommand{\ReLU}{\operatorname{ReLU}}
\newcommand{\CVaR}{\operatorname{CVaR}}
\newcommand{\VaR}{\operatorname{VaR}}
\newcommand{\SQ}{\operatorname{SQ}}
\newcommand{\Q}{\operatorname{Q}}

\title{Chance-Constrained Optimization via Bilevel Quantile Reformulation and Hypergradients}

% Anonymous for now; fill for camera-ready.
% \author[1]{...}
% \affil[1]{...}

\begin{document}
\maketitle

\begin{abstract}
Modern AI systems must optimize performance while controlling the probability of rare but consequential failures induced by stochasticity (noise, corruptions, distribution shift, or uncertain dynamics). Chance-constrained optimization provides a direct formalism for this goal, but is computationally challenging due to non-smooth quantile constraints and poor scaling with dimension and sample size.
We present a bilevel approach that reduces the chance constraint to a scalar inner problem via a superquantile variational characterization, then differentiates through the inner argmin using hypergradients. We combine a smooth strongly-convex surrogate for the inner objective with an augmented-Lagrangian-style penalty to obtain an end-to-end differentiable outer objective suitable for large-scale gradient methods.
We provide conditions under which the solution map is differentiable, derive an efficient gradient formula, and discuss approximation and optimization trade-offs. Experiments on standard synthetic benchmarks compare our method to TACO and conformal prediction based baselines and evaluate objective quality, reliability (empirical coverage), and runtime, including an additional coverage F-score metric.
\end{abstract}

\section{Introduction}\label{sec:intro}
\plan{
\begin{itemize}
  \item Motivate reliability in modern AI: robotics, perception under corruptions, safe decision-making.
  \item State the core problem: optimize utility subject to failure probability $\le \delta$.
  \item Explain why this is hard: quantiles, non-smoothness, sample-based combinatorics, scaling with $d$ and $M$.
  \item High-level idea: quantile $\rightarrow$ scalar bilevel $\rightarrow$ implicit differentiation $\rightarrow$ SGD-friendly objective.
  \item Contributions list (bullets) and preview experiments.
\end{itemize}
}

We study chance-constrained optimization (CCO), a framework for optimizing an objective while ensuring that constraint violations occur with probability at most $\delta$.
Such constraints naturally arise in modern AI systems that operate under uncertainty: a perception model must remain confident under random corruptions; a robot must avoid collisions under noisy sensing; and a policy must satisfy safety constraints under stochastic dynamics.

The main obstacle is that chance constraints are governed by quantiles of random variables, leading to non-smooth, typically non-convex feasible sets and optimization objectives that are difficult to differentiate and scale.
Sample-based encodings (e.g., mixed-integer formulations) and conformal-prediction-based methods can be effective on small problems but may struggle as the decision dimension or sample size grows.

\paragraph{Contributions.}
\begin{itemize}
  \item We formulate CCO as a bilevel problem using a quantile characterization and a superquantile (CVaR) variational representation.
  \item We introduce a smooth, strongly convex surrogate that yields a differentiable solution map and enables hypergradient computation.
  \item We combine this with an augmented-Lagrangian-style penalty to obtain a differentiable outer objective and a practical stochastic optimization algorithm.
  \item We benchmark against TACO and conformal prediction based solvers on synthetic problems and report objective value, empirical coverage, runtime, and a coverage F-score.
\end{itemize}

\section{Problem Setup}\label{sec:setup}
\subsection{Chance-Constrained Optimization}
Let $\Xc\subseteq \mathbb{R}^d$ be a nonempty closed convex set, $Z\sim P_Z$ a random vector on $\mathbb{R}^m$, $f:\mathbb{R}^d\to\mathbb{R}$ an objective, and $g:\mathbb{R}^d\times\mathbb{R}^m\to\mathbb{R}$ a (random) constraint function.
We study the chance-constrained optimization problem
\begin{align}
  \min_{x\in \Xc}\; & f(x)\label{eq:cco}\\
  \text{s.t.}\; & \mathbb{P}\big(g(x,Z)\le 0\big)\ge 1-\delta,\nonumber
\end{align}
where $\delta\in(0,1)$ controls the maximum acceptable probability of violation.
Equivalently, the violation probability satisfies $\mathbb{P}(g(x,Z)>0)\le \delta$.

\paragraph{Convexity.}
Throughout, we assume $f$ is convex and (sub)gradient-accessible.
Moreover, for every $z$, the map $x\mapsto g(x,z)$ is convex (and typically differentiable), and $\Xc$ is convex so that projection or proximal steps are well-defined.
Importantly, even under these convexity assumptions, the feasible set
\[
  \{x\in\Xc:\mathbb{P}(g(x,Z)\le 0)\ge 1-\delta\}
\]
is generally \emph{not} convex because it is defined through a quantile of the random variable $g(x,Z)$.
This non-smooth, nonconvex geometry is a major source of algorithmic difficulty.

\subsection{Access Model and Sampling}
We do not assume closed-form knowledge of $P_Z$.
Instead, we assume access to i.i.d.\ samples $z^{(1)},\dots,z^{(M)}\sim P_Z$, which is the standard setting in data-driven optimization and learning.
In experiments, we distinguish:
\begin{itemize}
  \item \emph{Training samples} used internally by an algorithm (e.g., to approximate expectations or quantiles).
  \item \emph{Evaluation samples} used only to estimate out-of-sample reliability.
\end{itemize}

\subsection{Empirical Coverage (Reliability Metric)}
For a candidate solution $x$, we measure reliability via the \emph{empirical coverage} computed on fresh i.i.d.\ samples $\tilde z^{(1)},\dots,\tilde z^{(N)}\sim P_Z$:
\begin{equation}
  \widehat{\mathrm{EC}}_N(x)
  \;:=\;
  \frac{1}{N}\sum_{i=1}^N \mathbf{1}\!\left\{g\!\left(x,\tilde z^{(i)}\right)\le 0\right\}.\label{eq:ec}
\end{equation}
The chance constraint in \eqref{eq:cco} targets $\mathrm{EC}(x)=\mathbb{P}(g(x,Z)\le 0)\ge 1-\delta$, while $\widehat{\mathrm{EC}}_N(x)$ provides an unbiased Monte Carlo estimate of $\mathrm{EC}(x)$.

\paragraph{Remark on extensions.}
For clarity, we focus on a single chance constraint in \eqref{eq:cco}; extensions to multiple constraints are discussed in Section~\ref{sec:extensions}.

\section{Background: Quantiles and Superquantiles}\label{sec:background}
Chance constraints are governed by quantiles of the random variable $g(x,Z)$, and it is useful to expose this structure explicitly.
This section derives a bilevel formulation in which the chance constraint is enforced through a scalar inner optimization problem.
We then explain why the resulting solution map $x\mapsto s^\star(x)$ can be unique but non-differentiable, motivating the relaxation in Section~\ref{sec:method}.

\subsection{Quantile reformulation}
For a scalar random variable $Y$ with CDF $F_Y$, define its $(1-\delta)$-quantile (value-at-risk) as
\begin{equation}
  \Q_{1-\delta}(Y)\;:=\;\inf\{t\in\mathbb{R}: F_Y(t)\ge 1-\delta\}.\label{eq:quantile-def}
\end{equation}
Applying this to $Y=g(x,Z)$, the chance constraint in \eqref{eq:cco} is equivalent to
\begin{equation}
  \mathbb{P}\big(g(x,Z)\le 0\big)\ge 1-\delta
  \quad\Longleftrightarrow\quad
  \Q_{1-\delta}\big(g(x,Z)\big)\le 0.\label{eq:quantile}
\end{equation}
Thus, feasibility is controlled by a single scalar summary of the distribution of $g(x,Z)$.

\subsection{Superquantiles and a scalar inner problem}
The \emph{superquantile} (conditional value-at-risk) at level $1-\delta$ is
\begin{equation}
  \SQ_{1-\delta}(Y)\;:=\;\mathbb{E}\!\left[Y\mid Y\ge \Q_{1-\delta}(Y)\right].\label{eq:sq-def}
\end{equation}
A key fact is that superquantiles admit a variational representation as the minimum of a convex function in a scalar auxiliary variable $s$ \citep{rockafellaruryasev2000cvar}:
\begin{equation}
  \SQ_{1-\delta}(Y)=\min_{s\in\mathbb{R}} \left\{ s + \frac{1}{\delta}\,\mathbb{E}\big[(Y-s)_+\big] \right\}.\label{eq:cvar-var}
\end{equation}
The minimizer set in \eqref{eq:cvar-var} is closely related to the quantile $\Q_{1-\delta}(Y)$; in particular, if $Y$ has a continuous distribution (no atom at the quantile), then the minimizer is unique and equals $\Q_{1-\delta}(Y)$ \citep{rockafellaruryasev2000cvar}.

Specializing to $Y=g(x,Z)$ motivates the scalar inner objective
\begin{equation}
  G_\delta(x,s)\;:=\; s + \frac{1}{\delta}\,\mathbb{E}_Z\!\left[(g(x,Z)-s)_+\right].\label{eq:G}
\end{equation}
Define the associated solution map
\begin{equation}
  s^\star(x)\;\in\;\arg\min_{s\in\mathbb{R}} G_\delta(x,s).\label{eq:sstar-def}
\end{equation}
When the minimizer is unique, we write $s^\star(x)$ for the unique scalar minimizer.
Under the continuity condition above, $s^\star(x)$ coincides with the $(1-\delta)$-quantile of $g(x,Z)$, hence the chance constraint \eqref{eq:quantile} is equivalent to $s^\star(x)\le 0$.
This yields the bilevel viewpoint of CCO developed in \citet{laguelmalickackooij2021bilevel}.

\begin{lemma}[Uniqueness of the inner minimizer]\label{lem:unique-sstar}
Fix $x\in\Xc$ and define $Y:=g(x,Z)$.
Assume $Z$ is absolutely continuous with respect to Lebesgue measure and $g(x,\cdot)$ is differentiable.
If $Y$ has a continuous distribution (equivalently, $\mathbb{P}(Y=t)=0$ for all $t\in\mathbb{R}$), then $\arg\min_{s\in\mathbb{R}} G_\delta(x,s)$ is a singleton and the unique minimizer satisfies
\[
  s^\star(x)=\Q_{1-\delta}(Y)=\Q_{1-\delta}(g(x,Z)).
\]
\end{lemma}
\begin{proof}
For fixed $x$, the function $s\mapsto G_\delta(x,s)$ is convex.
Let $F_Y$ be the CDF of $Y$.
Standard subgradient calculus for the hinge yields
\[
  \partial_s G_\delta(x,s)
  \;=\;
  \left[\,1-\frac{1}{\delta}\mathbb{P}(Y>s)\;,\;1-\frac{1}{\delta}\mathbb{P}(Y\ge s)\,\right].
\]
Thus $0\in \partial_s G_\delta(x,s)$ if and only if $\mathbb{P}(Y>s)\le \delta \le \mathbb{P}(Y\ge s)$, i.e., $s$ belongs to the $(1-\delta)$-quantile set of $Y$.
If $Y$ has no atom at the quantile (in particular, if $Y$ is continuous), then this set reduces to a single point, so the minimizer is unique and equals $\Q_{1-\delta}(Y)$.
\end{proof}

\subsection{Uniqueness vs.\ differentiability of the quantile map}
For each fixed $x$, the function $s\mapsto G_\delta(x,s)$ is convex and continuously differentiable.
Moreover, if $g(x,Z)$ has a continuous distribution, then the first-order condition $\partial_s G_\delta(x,s)=0$ has a unique solution, implying that $s^\star(x)$ is unique.

However, \emph{uniqueness does not imply differentiability} of the map $x\mapsto s^\star(x)$.
Intuitively, the kink in $(u)_+$ causes the derivatives of $G_\delta$ to involve indicator functions that change discontinuously as $x$ varies.
Formally, differentiating through the optimality condition $\partial_s G_\delta(x,s^\star(x))=0$ would require second derivatives in $s$ (or an implicit function theorem argument), but $G_\delta(x,s)$ is not twice differentiable in $s$ because of the hinge.
As a consequence, standard differentiability results for argmin maps (e.g., via Danskin-type theorems requiring $C^2$ smoothness and strong convexity in the inner variable \citep{danskin1967theory}) do not apply directly.
Section~\ref{sec:method} addresses this by introducing a smooth, strongly convex surrogate whose minimizer is differentiable.

\section{Method: Bilevel Augmented Hypergradient Optimization}\label{sec:method}
We now describe our bilevel augmented approach.
Starting from the bilevel reformulation in \eqref{eq:sstar-def}, our goal is to optimize $f(x)$ while driving the quantile proxy $s^\star(x)$ below zero.
The main challenge is that $s^\star(x)$ can be unique but non-differentiable (Section~\ref{sec:background}), which obstructs direct gradient-based optimization.
We address this by (i) smoothing and strongly convexifying the inner problem to obtain a differentiable surrogate $\widehat s^\star(x)$, and (ii) penalizing constraint violation through a differentiable augmented objective.

\subsection{From chance feasibility to a bilevel constraint}
Under the conditions described in Section~\ref{sec:background}, the chance constraint is equivalent to $s^\star(x)\le 0$, where
\[
  s^\star(x)\in\arg\min_{s\in\mathbb{R}} G_\delta(x,s),
  \qquad
  G_\delta(x,s)= s + \frac{1}{\delta}\,\mathbb{E}\!\left[(g(x,Z)-s)_+\right].
\]
This yields the constrained outer problem
\begin{equation}
  \min_{x\in\Xc}\; f(x)\qquad \text{s.t.}\qquad s^\star(x)\le 0.\label{eq:outer-constrained}
\end{equation}

\subsection{A differentiable augmented objective}
We enforce $s^\star(x)\le 0$ via a standard Lagrangian relaxation. The constrained problem \eqref{eq:outer-constrained} can be written as
\begin{equation}
  \min_{x\in\Xc}\;\max_{\lambda\ge 0}\; f(x) + \lambda s^\star(x).\label{eq:minmax}
\end{equation}
The max in \eqref{eq:minmax} is non-smooth in $x$.
We therefore introduce a quadratic regularization on the dual variable (an augmented-Lagrangian-style smoothing):
\begin{equation}
  \min_{x\in\Xc}\;\max_{\lambda\ge 0}\; f(x) + \lambda s^\star(x) - \frac{\mu}{2}\lambda^2,\label{eq:aug-lag}
\end{equation}
where $\mu>0$ controls the strength of the regularization.
Maximizing over $\lambda$ yields the equivalent penalized objective
\begin{equation}
  F_\mu(x)\;:=\; f(x) + \frac{([s^\star(x)]_+)^2}{2\mu}.\label{eq:Fmu}
\end{equation}
When $s^\star(x)\le 0$, the penalty term vanishes; when $s^\star(x)>0$, it grows quadratically.

\subsection{Relaxing the inner problem: smoothness and strong convexity}\label{sec:inner-relax}
To optimize \eqref{eq:Fmu} with gradients, we need to differentiate through the map $x\mapsto s^\star(x)$.
As discussed in Section~\ref{sec:background}, $s^\star(x)$ is generally not differentiable because $(u)_+$ is not $C^2$.
We therefore replace the hinge with a $C^2$ smoothing and add a strongly convex term in $s$.

\paragraph{Smooth approximation of the hinge.}
We use a twice continuously differentiable approximation $h_\theta:\mathbb{R}\to\mathbb{R}$ to $(\cdot)_+$ with smoothing parameter $\theta>0$:
\begin{equation}
  h_\theta(u)=
  \begin{cases}
    -\dfrac{u^4}{16\theta^3} + \dfrac{3u^2}{8\theta} + \dfrac{u}{2} + \dfrac{3\theta}{16}, & |u|\le \theta,\\[6pt]
    \max\{u,0\}, & \text{otherwise.}
  \end{cases}\label{eq:htheta}
\end{equation}

\paragraph{Strongly convex surrogate.}
Define the surrogate inner objective
\begin{equation}
  \widehat G_{\delta,\theta,\varepsilon}(x,s)
  \;:=\;
  s + \frac{1}{\delta}\,\mathbb{E}_Z\!\left[h_\theta\!\left(g(x,Z)-s\right)\right]
  + \frac{\varepsilon}{2}s^2,\label{eq:Ghat}
\end{equation}
with $\varepsilon>0$.
For each $x$, the function $s\mapsto \widehat G_{\delta,\theta,\varepsilon}(x,s)$ is $\varepsilon$-strongly convex and $C^2$, hence it has a unique minimizer
\begin{equation}
  \widehat s^\star(x)\;:=\;\arg\min_{s\in\mathbb{R}} \widehat G_{\delta,\theta,\varepsilon}(x,s).\label{eq:shat-star}
\end{equation}

\subsection{Hypergradients via implicit differentiation}
Because the inner problem \eqref{eq:shat-star} is strongly convex and smooth in $s$, the solution map $x\mapsto \widehat s^\star(x)$ is differentiable under mild regularity assumptions on $g$.
Let $\phi(x,s):=\partial_s \widehat G_{\delta,\theta,\varepsilon}(x,s)$.
The optimality condition $\phi(x,\widehat s^\star(x))=0$ and the fact that $\partial_s\phi(x,s)=\partial^2_{ss}\widehat G_{\delta,\theta,\varepsilon}(x,s)\ge \varepsilon$ imply (by the implicit function theorem) that $\widehat s^\star(x)$ is $C^1$ and
\begin{equation}
  \nabla_x \widehat s^\star(x)
  \;=\;
  - \frac{\partial^2_{xs}\widehat G_{\delta,\theta,\varepsilon}(x,s)}{\partial^2_{ss}\widehat G_{\delta,\theta,\varepsilon}(x,s)}
  \bigg|_{s=\widehat s^\star(x)}.\label{eq:hypergrad}
\end{equation}
We then optimize the differentiable surrogate of \eqref{eq:Fmu}:
\begin{equation}
  \widehat F_\mu(x)\;:=\; f(x) + \frac{([\widehat s^\star(x)]_+)^2}{2\mu}.\label{eq:Fhat}
\end{equation}

\subsection{Sample-based implementation}
In practice, we approximate expectations by sample averages.
Given a minibatch $\{z^{(j)}\}_{j=1}^M$, we use the plug-in estimator
\begin{equation}
  \widehat G_{\delta,\theta,\varepsilon}(x,s;\{z^{(j)}\}_{j=1}^M)
  \;:=\;
  s + \frac{1}{\delta M}\sum_{j=1}^M h_\theta\!\left(g(x,z^{(j)})-s\right) + \frac{\varepsilon}{2}s^2.\label{eq:Ghat-plugin}
\end{equation}
The inner minimizer $\widehat s^\star(x)$ can be obtained efficiently because it is one-dimensional and the derivative $\partial_s \widehat G$ is strictly increasing in $s$; we use a robust bisection/Newton procedure.
The outer step applies a gradient-based update (e.g., SGD/Adam) to $\widehat F_\mu(x)$, optionally projecting back onto $\Xc$.
We may decrease $\mu$ over iterations to sharpen feasibility while keeping gradients stable.

\subsection{Algorithm (two-time-scale view)}
\begin{figure}[!t]
  \centering
  \fbox{\begin{minipage}{0.97\linewidth}
  \textbf{Bilevel Augmented Hypergradient Method (CCO-BLO).}
  \begin{enumerate}
    \item Initialize $x_0\in\Xc$ and $s_0\in\mathbb{R}$. Choose $\delta,\theta,\varepsilon,\{\mu_k\}$.
    \item For $k=0,1,\dots,K-1$:
      \begin{enumerate}
        \item Sample minibatch $\{z^{(j)}\}_{j=1}^M\sim P_Z$.
        \item (Inner) Compute $s_{k+1}\approx \widehat s^\star(x_k)$ by minimizing \eqref{eq:Ghat-plugin} in $s$.
        \item (Outer) Form $\nabla_x \widehat F_{\mu_k}(x_k)$ using \eqref{eq:hypergrad} and update $x_{k+1}$ with a gradient step (and project onto $\Xc$ if needed).
      \end{enumerate}
  \end{enumerate}
  \end{minipage}}
  \caption{High-level procedure; see Section~\ref{sec:inner-relax} for definitions.}
  \label{fig:algorithm}
\end{figure}

\subsection{Extensions: multiple constraints}\label{sec:extensions}
For separate chance constraints $\mathbb{P}(g_i(x,Z)\le 0)\ge 1-\delta_i$, $i=1,\dots,r$, we introduce one scalar $s_i$ per constraint and apply the same construction to each inner objective $\widehat G_{\delta_i,i}(x,s_i)$.
The inner step decomposes into $r$ independent one-dimensional solves, and the outer penalty can be taken as $\sum_{i=1}^r ([\widehat s_i^\star(x)]_+)^2/(2\mu_i)$ (or with a shared $\mu$).
This keeps the inner dimensionality small even when the decision variable $x$ is high-dimensional.

\section{Theory (What We Will Prove)}\label{sec:theory}
\plan{
\begin{itemize}
  \item State assumptions: $g$ differentiable in $x$, mild continuity in $z$, bounded moments; $\varepsilon>0$.
  \item Prove inner uniqueness and differentiability of $\widehat s^\star(x)$ (implicit function theorem).
  \item Provide bounds relating $\widehat s^\star$ to the original quantile and discuss $\theta,\varepsilon\to 0$.
  \item Provide a convergence statement for stochastic optimization of $\widehat F_\mu$ (stationary point) under standard stepsizes.
  \item Discuss extensions: multiple constraints (vector $s$), joint chance (risk allocation).
\end{itemize}
}

This section will formalize conditions under which the hypergradient \eqref{eq:hypergrad} is valid and stable.
We will also characterize the approximation error introduced by smoothing and strong convexification and provide optimization guarantees for the resulting stochastic bilevel procedure.

\paragraph{Core proof skeleton.}
\begin{itemize}
  \item Show $\partial^2_{ss}\widehat G_{\delta,\theta,\varepsilon}(x,s)\ge \varepsilon>0$ for all $x,s$, implying unique $\widehat s^\star(x)$.
  \item Apply the implicit function theorem to $\partial_s \widehat G_{\delta,\theta,\varepsilon}(x,\widehat s^\star(x))=0$ to obtain $C^1$ smoothness and \eqref{eq:hypergrad}.
  \item Provide a bound of the form $|\widehat s^\star(x)-s^\star(x)|\le c_1\theta + c_2\varepsilon$ under regularity conditions (to be specified).
\end{itemize}

\section{Experiments}\label{sec:experiments}
\plan{
\begin{itemize}
  \item Synthetic benchmarks aligned with the deck:
    (i) CPP univariate exponential example; (ii) high-dimensional Gaussian linear constraint with varying $d$ and $M$.
  \item Compare against baselines: TACO; CPP-MIP/CPP-KKT/SA/SAA.
  \item Metrics: objective, empirical coverage, runtime; plus coverage F-score (closeness to target $1-\delta$).
  \item Sweeps: dimension scaling ($d$), sample scaling ($M$), and sensitivity to $\delta$.
  \item Report failure modes: CPP timeouts, instability, dimension/sample sensitivity.
\end{itemize}
}

\subsection{Setup}
We follow standard synthetic benchmarks used in prior work \citep{zhao2024conformal}.
We report the objective value and empirical coverage (estimated on fresh samples) along with wall-clock time.
To summarize reliability in a single scalar, we additionally report a coverage F-score:
\begin{equation}
  F_{\mathrm{EC}}(\mathrm{EC}, 1-\delta)
  \;=\;
  \frac{2\min\{\mathrm{EC}, 1-\delta\}}{\mathrm{EC} + (1-\delta)}\in[0,1],\label{eq:fscore}
\end{equation}
which equals $1$ if and only if the empirical coverage equals the target.

\subsection{Baselines}
We compare against (i) TACO \citep{laguelmalickackooij2021bilevel} and (ii) conformal predictive programming methods (MIP and KKT encodings, plus SA/SAA variants) \citep{zhao2024conformal}.

\subsection{Results (placeholders)}
\todo{Insert tables/figures from runs/ once sweeps finish on the server.}

\begin{table}[!t]
  \centering
  \caption{Synthetic Benchmark Summary (placeholder).}\label{tab:main-results}
  \begin{tabular}{lrrrr}
    \toprule
    Method & Obj.\ $\downarrow$ & EC $\uparrow$ & $F_{\mathrm{EC}}\uparrow$ & Time $\downarrow$\\
    \midrule
    Ours & -- & -- & -- & --\\
    TACO & -- & -- & -- & --\\
    CPP-MIP & -- & -- & -- & --\\
    CPP-KKT & -- & -- & -- & --\\
    \bottomrule
  \end{tabular}
\end{table}

\section{Discussion and Limitations}\label{sec:discussion}
\plan{
\begin{itemize}
  \item When the method is attractive: high-dimensional $x$ with a small number of chance constraints.
  \item Sensitivity knobs: $\theta,\varepsilon,\mu$ schedule; batch size $M$; inner solve accuracy.
  \item Limitations: bias from smoothing; local minima in nonconvex $f$; joint chance constraints; heavy tails.
  \item Practical guidance: penalty schedule; robust inner solve (bisection); reporting reliability metrics.
\end{itemize}
}

\section{Conclusion}\label{sec:conclusion}
\plan{
\begin{itemize}
  \item Restate core idea: scalar inner problem + hypergradients + penalty.
  \item Summarize empirical takeaways.
  \item Roadmap: modern AI experiment (robust classification under random corruptions), multi-constraint extension.
\end{itemize}
}

\bibliography{references}

\end{document}
